% vim: et
% common packages, fonts, styles and commands
% unlocalised content
\documentclass[a4paper,9pt]{article}
\usepackage{fontspec}
\usepackage[a4paper,lmargin=2cm,rmargin=2cm,bottom=2cm,top=2cm]{geometry}
\usepackage{multicol}
\usepackage{titling}
\usepackage{fancyhdr}
\usepackage{paralist}
\usepackage{gensymb}

\setmainfont{Roboto Slab Light}

% section without title
\newcommand{\nosection}[1]{%
  \refstepcounter{section}%
  \addcontentsline{toc}{section}{\protect\numberline{\thesection}#1}%
  \markright{#1}}

% title formatting
\newfontfamily\headingfont[]{Oswald}
\newfontfamily\constfont[]{Roboto}
\newcommand{\constt}[1]{{\constfont\selectfont{\textit{#1}}}}
\renewcommand{\maketitlehooka}{\headingfont}
\setlength{\droptitle}{-5em}

% paragraph spacing
\setlength{\parindent}{0em}
\setlength{\parskip}{0.5em}


\usepackage{polyglossia}
\setdefaultlanguage{slovak}
\usepackage{background}
\usetikzlibrary{calc}
\usepackage{fixltx2e}
\usepackage{unicode-math}
\usepackage{fancybox}
\usepackage{amsmath}
\usepackage{graphicx}
\usepackage{tabularx}

% lines for cutting and bending for folded book like manual
\backgroundsetup{
color=black,
scale=1,
opacity=1,
angle=0,
contents={\tikz\draw[line width=0.1mm, gray]
    (0mm, 0mm) -- (0mm, 297mm)
    (180mm, 0mm) -- (180mm, 297mm)
    [dashed]
    (60mm, 0mm) -- (60mm, 297mm)
    (120mm, 0mm) -- (120mm, 297mm);
  }
}

\newcommand{\makefulltitle}{Ako používať Logaritmické pravítko }
% localisation strings used in the templates
\newcommand{\tabstockupper}{Horná polovica Pevnej časti }
\newcommand{\tabslidefront}{Predná strana Posuvnej časti }
\newcommand{\tabslideback}{Zadná strana Posuvnej časti }
\newcommand{\tabstocklower}{Spodná polovica Pevnej časti }
% custom page number
\fancypagestyle{plain}{% title page
  \renewcommand{\headrulewidth}{0pt}%
  \fancyhf{}%
  \fancyfoot[L]{\headingfont{http://wheel.creat.io/sr}}%
  \fancyfoot[R]{\headingfont{\makefulltitle \textbf{\thepage}}}%
  \renewcommand{\headrulewidth}{0pt}
}
\pagestyle{fancy}% other pages
\fancyhf{}
\lfoot{\headingfont{http://wheel.creat.io/sr}}
\rfoot{\headingfont{\makefulltitle \textbf{\thepage}}}
\renewcommand{\headrulewidth}{0pt}


% common part of how-to perex
\newcommand{\makeperex}{Logaritmické pravítko je analógový počítač - zariadenie, ktoré dokáže počítať väčšinu základných a niektoré pokročilé matematické operácie. Nie je síce také presné a nepoužíva sa tak jednoducho ako kalkulačka, ale elegantne demonštruje matematické zákony. Je to historický artefakt, ktorý bol ešte pred päťdesiatimi rokmi bežne používaný v technických disciplínach, inžinieri a astronauti ním počas vesmírnych misií dokonca počítali trajektórie a letové parametre. }
% custom page number
\fancypagestyle{plain}{% title page
  \renewcommand{\headrulewidth}{0pt}%
  \fancyhf{}%
  \fancyfoot[L]{\headingfont{http://wheel.creat.io/sr}}%
  \fancyfoot[R]{\headingfont{\makefulltitle \textbf{\thepage}}}%
  \renewcommand{\headrulewidth}{0pt}
}
\pagestyle{fancy}% other pages
\fancyhf{}
\lfoot{\headingfont{http://wheel.creat.io/sr}}
\rfoot{\headingfont{\makefulltitle \textbf{\thepage}}}
\renewcommand{\headrulewidth}{0pt}





\title{\fontsize{60}{60}\selectfont LOGARITMICKÉ PRAVÍTKO}
\preauthor{}\postauthor{}\author{}
\predate{}\postdate{}\date{}
\begin{document}

  \begin{center}
    \headingfont\fontsize{28}{28}\selectfont AKO POUŽÍVAŤ LOGARITMICKÉ PRAVÍTKO
  \end{center}

%  {\let\newpage\relax\maketitle}% no new line before title
  \nosection{}
  \large\textbf{\makeperex Tento návod popisuje ako používať papierové Logaritmické pravítko \modelname.}

  \begin{multicols*}{3}
  \normalsize{
  Logaritmické pravítko sa skladá z:
    \begin{inparaenum}[a\upshape)]
      \item Pevnej časti;
      \item Posuvnej časti; a
      \item Bežca.
    \end{inparaenum}

  Na prednej strane Pevnej časti a oboch stranách Posuvnej časti sú zobrazené logaritmické stupnice, označené veľkými písmenami:

  % table of stock and slide scales
\renewcommand{\arraystretch}{1.2}
\begin{center}
  \begin{tabular}{ |c|c|}
    \hline
    \multicolumn{2}{|c|}{\tabstockupper} \\
    \hline
      L & \(log x\) \\
    \hline
      P & \(\sqrt{1 - x^2}\) \\
    \hline
      K & \(x^3\) \\
    \hline
      A & \(x^2\) \\
    \hline
    \hline
    \multicolumn{2}{|c|}{\tabslidefront} \\
    \hline
      B & \(x^2\) \\
    \hline
      CI & \(1/x\) \\
    \hline
      C & \(x\) \\
    \hline
    \hline
    \multicolumn{2}{|c|}{\tabslideback} \\
    \hline
      LL1 & \(e^{0.01x}\) \\
    \hline
      LL2 & \(e^{0.1x}\) \\
    \hline
      LL3 & \(e^{1x}\) \\
    \hline
    \hline
    \multicolumn{2}{|c|}{\tabstocklower} \\
    \hline
      D & \(x\) \\
    \hline
      S & \(sin x\) \\
    \hline
      T & \(tan x\) \\
    \hline
      ST & \(sin x \approx tan x \) \\
    \hline
  \end{tabular}
\end{center}


  Čiary na stupniciach sa nazývajú \textbf{indexy} alebo rysky. Číselné značky na stupnici označujú významné indexy.

  Na zadnej strane Pevnej časti sú tabuľky matematických a fyzikálnych konštánt a ďalšia užitočná grafika\footnote{Metrické a palcové stupnice na zadnej strane Pevnej časti a prednej strane Bežca nemusia merať presne, pretože teplota a vlhkosť spôsobené lepením menia mechanické vlastnosti a rozmery papiera.}, ale tieto sa priamo nepoužívajú v matematických operáciách. Posuvná časť  obsahuje prevodné mierky pre imperiálne a metrické (SI) jednotky a na vnútorných stranách prevodné stupnice medzi číselnými sústavami a mierami uhlov.
a
  Posuvná časť sa pohybuje vo vnútri Pevnej časti a Bežec po jej povrchu. Matematické operácie sa vykonávajú vzájomným posúvaním častí, zarovnaním indexov stupníc a odčítaním hodnôt. Niektoré operácie vyžadujú drobnú konverziu desatinných miest.

  \centerline{\includegraphics{ops_sk}}
  \vspace*{3mm}

  \makeop{2mm}{{x}\cdot{y}}
  \textbf{Jednoduché násobenie}

\constt{1.2{\char"00D7}2.3:}
Posuň index 1 stupnice C na pozíciu indexu 1.2 stupnice D.
Posuň Bežec na index 2.3 stupnice D.
Hodnota, ktorú Bežec ukazuje na stupnici C je výsledok: \constt{2.76}.

\footnotesize Použi stupnice A a B pre násobenie a delenie väčších čísel.\normalsize

  \textbf{Násobenie s prechodom}
\constt{2.4{\char"00D7}4.6}
Posuň index 10 stupnice C na pozíciu indexu 2.4 stupnice D.
Posuň Bežec na index 4.6 stupnice C.
Hodnota, ktorú Bežec ukazuje na stupnici D je 1.105.
Z pamäti vypočítaj, že súčin približných hodnôt 2.5{\char"00D7}5 je 12.5.
Výsledok je teda väčší ako 10.
Úpravou desatinného miesta získaš výsledok: \constt{11.05}. 

  \makeop{2mm}{{x}/{y}}
  \textbf{Jednoduché delenie}

\constt{4.6/7.7:}
Posuň Bežec na index 4.6 stupnice D.
Posuň index 7.7 stupnice C na Bežec.
Posuň Bežec buď na index 1 alebo index 10 stupnice C, podľa toho, ktorý je bližšie. V tomto prípade index 10.
Bežec je teraz na indexe 5.97 na stupnici D. Správna odpoveď je približne 4 / 8 = 0.5. Úpravou desatinného miesta získaš výsledok: \constt{0.597}.

  \makeop{2mm}{{1}/{x}}
  \textbf{Recipročná hodnota}

\footnotesize Stupnica CI je inverzná stupnica, hodnoty indexov stúpajú zprava doľava. \normalsize

\constt{1/7.6:}
Posuň Bežec na index 7.8 stupnice CI.
Bežec je teraz na indexe 1.132 na stupnici C.
Správna odpoveď je približne 1 / 10 = 0.1. Úpravou desatinného miesta získaš výsledok: \constt{0.1132}.

  \makeop{3mm}{x^2}
  \textbf{Druhá mocnina}

\constt{4.2\textsuperscript{2}:}
Posuň Bežec na index 4.2 stupnice C.
Bežec ukazuje výsledok \constt{17.6} na stupnici B.

  \makeop{3mm}{\sqrt{x}}
  \textbf{Druhá odmocnina}

\footnotesize Pre počítanie druhej odmociny čísel s nepárnym počnom číslic použi rozsah [1, 10] stupnice B. Pre druhú mocninu čísel s párnym počtom číslic použi rozsah [10, 100] stupnice B. \normalsize

\constt{4400\textsuperscript{0.5}:}
Posuň Bežec na index 44 stupnice B (4400 má párny počet číslic).
Bežec je na indexe 6.65 na strupnici C. Správna odpoveď je približne 70, keďźe 70\textsuperscript{2} = 4900. Úpravout desatinného miesta získaš výsledok: \constt{66.5}.

  \makeop{3mm}{x^3}
  \textbf{Tretia mocnina}

\constt{4.8\textsuperscript{3}:}
Posuň Bežec na index 4.8 stupnice D.
Bežec ukazuje výsledok \constt{110} na stupnici K.

  \makeop{3mm}{\sqrt[\leftroot{5}\uproot{2}3]{x}}
  \textbf{Tretia odmocnina}

\footnotesize Rozsah [1, 10] stupnice K sa používa pre tretie odmocniny čísel s jednou číslicou, rozsah [10, 100] pre čísla s dvoma číslicami, [100, 1000] pre trojmiestne čísla. Čísla so štyrmi číslicami opäť prvý rozsah, päť číslic druhý rozsah, atď. \normalsize

\constt{4400\textsuperscript{0.33}:}
Posuň Bežec na index 4.4 stupnice K (4400 má 4 číslice).
Bežec je na indexe 1.655 na stupnici D.
Správna odpoveď je medzi 10 a 20 pretože 10\textsuperscript{3} = 1000 a 20\textsuperscript{3} = 8000. Úpravou desatinného miesta získaš výsledok: \constt{16.55}.

  \makeop{3mm}{{x}^{10}}
  \textbf{Desiata mocnina}

\footnotesize Mocnenie na mocniny čísla 10 a všeobecné mocniny (vrátane odmocnín) sa počítajú na stupniciach LL, ktoré sú umiestnené na opačnej strane Posuvnej časti. \normalsize

\constt{1.36\textsuperscript{10}:}
Posuň Bežec na index 1.36 stupnice LL2.
Bežec sa nachádza na indexe \constt{21.5} na stupnici LL3, čo je výsledok operácie.

  \makeop{3mm}{\sqrt[\leftroot{6}\uproot{2}10]{x}}
  \textbf{Desiata odmocnina}

\constt{5\textsuperscript{0.1}:}
Posuň Bežec na index 5 stupnice LL3.
Bežec sa nachádza na indexe \constt{1.175} na stupnici LL2, čo je výsledok operácie.

  \makeop{3mm}{{x}^{100}}
  \textbf{Stá mocnina}

\constt{1.02\textsuperscript{100}:}
Posuň Bežec na index 1.02 stupnice LL1.
Bežec sa nachádza na indexe \constt{7.25} na stupnici LL3, čo je výsledok operácie.

  \makeop{3mm}{\sqrt[\leftroot{7}\uproot{2}100]{x}}
  \textbf{Stá odmocnina}

\constt{5\textsuperscript{0.01}:}
Posuň Bežec na index 5 stupnice LL3.
Bežec sa nachádza na indexe \constt{1.022} na stupnici LL1, čo je výsledok operácie.

  \makeop{3mm}{x^y}
  \textbf{Všeobecná mocnina}

\constt{9.1\textsuperscript{2.3}:}
Posuň index 9.1 stupnice LL3 na pozíciu indexu 1 stupnice D.
Posuň Bežec na index 2.3 stupnice D.
Pozícia Bežca na stupnici LL3 je výsledok: \constt{160.6}. Bežec je v skutočnosti v blízkosti 161, pretože stupnica stráca presnosť so stúpajúcou hodnotou. 

\constt{230\textsuperscript{0.45}:}
Posuň index 230 stupnice LL3 na pozíciu indexu 10 stupnice C.
Posuň Bežec na index 4.5 stupnice C.
Bežec sa nachádza na indexe výsledku na stupnici LL3: \constt{11.6}.

\constt{0.78\textsuperscript{3.4}:}
Posuň index 3.4 stupnice LL3 na pozíciu indexu 7.8 stupnice C.
Posuň Bežec na index 10 stupnice C.
Bežec sa nachádza na indexe výsledku na stupnici LL3: \constt{4.3}.

\constt{1.9\textsuperscript{2.5}:}
Posuň index 1.9 stupnice LL2 na pozíciu indexu 10 stupnice C.
Posuň Bežec na index 2.5 stupnice C.
Bežec sa nachádza na indexe výsledku na stupnici LL3: \constt{4.97}.

\constt{0.99\textsuperscript{560}:}
Posuň index 9.9 stupnice LL3 na pozíciu indexu 10 stupnice C.
Posuň Bežec na index 5.6 stupnice C.
Bežec sa nachádza na indexe 3.6 na stupnici LL3.
Pretože si použil/použila jednu stotinu exponentu, úpravou hodnoty o dve desatinné miesta získaš výsledok: \constt{0.036}.

  \makeop{3mm}{\sin{x}}
  \textbf{Sínus}

\footnotesize Použi stupnicu S pre výpočet sínusu uhlov medzi 5\textdegree a 90\textdegree, a stupnicu ST pre uhly medzi 30' a 6\textdegree. \normalsize

\constt{sin(33\textdegree):}
Posuň Bežec na 33 stupnice S.
Bežec sa nachádza na indexe 5.45 na stupnici C.
Správna odpoveď je medzi 0.1 a 1. Úpravou desatinného miesta získaš výsledok: \constt{0.545}.

  \makeop{3mm}{\cos{x}}
  \textbf{Kosínus}

\constt{cos(33\textdegree):}
Posuň Bežec na index 33 stupnici S.
Bežec sa nachádza na indexe výsledku \constt{0.83} na stupnici P (inverzná stupnica).

  \makeop{3mm}{\tan{x}}
  \textbf{Tangens}

\footnotesize Použi stupnicu T pre výpočet tangensu uhlov medzi 5\textdegree a 45\textdegree, a stupnicu ST pre uhly medzi 30' a 6\textdegree. \normalsize

\constt{tan(33\textdegree):}
Posuň Bežec na index 33 stupnice S.
Bežec sa nachádza na indexe 6.5 na stupnici C.
Správny výsledok je medzi 0.1 a 1. Úpravou desatinného miesta získaš výsledok: \constt{0.65}.

  \makeop{3mm}{\log_{10}{x}}
  \textbf{Dekadický logaritmus}

\constt{log\textsubscript{10}2.6:}
Posuň Bežec na index 2.6 stupnice D.
Bežec sa nachádza na indexe 4.14 na stupnici L.
Úpravou hodnoty o desatinné miesto získaš výsledok: \constt{0.414}.
  }
  \end{multicols*}
  
\end{document}
